\documentclass[12pt, a4paper]{report}
\usepackage[T1]{fontenc} % Proper < > | characters in code
\usepackage{lmodern}
\usepackage{graphicx} % For including images
\usepackage{amsmath} % For math mode commands and environments
\usepackage{amssymb} % For \checkmark
\usepackage{caption} % For listing and table captions
\usepackage{titlesec} % For customizing section titles
\usepackage{hyperref} % For clickable links in the document
\usepackage{booktabs} % For better tables
\usepackage{array} % For table column formatting
\usepackage{tabularx} % For tables that fill the text width
\usepackage{enumitem} % For customizing lists
\usepackage{xcolor} % For code colours
\usepackage{listings} % For source code listings
\usepackage{float} % For [H] table placement
%% ____Bibliography____%%
\usepackage[numbers,sort&compress]{natbib}
% Page margins
\usepackage[left=1.5in, right=1in, top=1in, bottom=1in]{geometry}

\hypersetup{hidelinks}

% A short report: numbered sections without chapters
\renewcommand{\thesection}{\arabic{section}}
\titleformat{\section}{\large\bfseries}{\thesection}{0.6em}{}
\titleformat{\subsection}{\normalsize\bfseries}{\thesubsection}{0.6em}{}
\titlespacing*{\section}{0pt}{10pt plus 2pt minus 2pt}{4pt plus 1pt}
\titlespacing*{\subsection}{0pt}{6pt plus 2pt minus 2pt}{2pt}

% Captions and references
\captionsetup{font=small, labelfont=bf, skip=3pt}
\renewcommand{\bibsection}{\section*{References}}
\renewcommand{\bibfont}{\footnotesize}
\setlength{\bibsep}{0pt}

% Source code
\lstdefinestyle{base}{
  basicstyle=\ttfamily\footnotesize,
  numbers=left, numberstyle=\tiny\color{gray}, numbersep=5pt,
  xleftmargin=14pt, framexleftmargin=14pt,
  frame=tb, framerule=0.4pt, rulecolor=\color{gray!70},
  backgroundcolor=\color{gray!6},
  keywordstyle=\color{blue!55!black}\bfseries,
  commentstyle=\color{green!35!black}\itshape,
  stringstyle=\color{red!55!black},
  columns=fullflexible, keepspaces=true, showstringspaces=false,
  upquote=true, captionpos=b, aboveskip=6pt, belowskip=2pt
}
\lstdefinestyle{hs}{style=base, language=Haskell}
\lstdefinestyle{py}{style=base, language=Python, numbers=none, xleftmargin=0pt,
  framexleftmargin=0pt, basicstyle=\ttfamily\fontsize{7.5}{8.1}\selectfont}
\lstdefinestyle{plain}{style=base, numbers=none, xleftmargin=0pt,
  framexleftmargin=0pt, basicstyle=\ttfamily\footnotesize}
\renewcommand{\lstlistingname}{Listing}

\newcommand{\code}[1]{\texttt{#1}}
\setlength{\emergencystretch}{2em} % Avoid overfull lines next to long code
\frenchspacing % No extra space after colons and full stops in code

\begin{document}


\thispagestyle{empty}

\begin{center}

\begin{center}
 \includegraphics[width=3cm,keepaspectratio=true]{uor_logo.jpg}
 % uor_logo.jpg: 236x331 pixel, 72dpi, 8.33x11.68 cm, bb=0 0 236 331
\end{center}

\vspace{1.5cm}
\begin{huge}
%%%%%%%%%%%%%%%%%%%%%%%%%%%%%%%%%%%%%%%%%%%%%%%%%%%%
% Project title
%%%%%%%%%%%%%%%%%%%%%%%%%%%%%%%%%%%%%%%%%%%%%%%%%%%%
A Type-Safe Arithmetic Expression Interpreter in Haskell
%%%%%%%%%%%%%%%%%%%%%%%%%%%%%%%%%%%%%%%%%%%%%%%%%%%%
\end{huge} \\
\vspace{1cm}

\begin{normalsize}
A project assignment report submitted to the
\end{normalsize}\\
\vspace{1cm}

\begin{large}
Department of Electrical and Information Engineering\\
Faculty of Engineering\\
University of Ruhuna\\
Sri Lanka
\end{large}\\

\vspace{1cm}

\begin{large}\textbf{EC 8206 Functional Programming}\end{large}\\

\vspace{1cm}
\begin{normalsize}by  \end{normalsize}
\vspace{1cm}

\begin{tabular}[h]{lll}
 %%%%%%%%%%%%%%%%%%%%%%%%%%%%%%%%%%%%%%%%%%%%%%%%%%%%%%%%%%
 % Name and Registration Number
 %%%%%%%%%%%%%%%%%%%%%%%%%%%%%%%%%%%%%%%%%%%%%%%%%%%%%%%%%%
 Achintha H.G.R. & - & EG/2021/4384
 %%%%%%%%%%%%%%%%%%%%%%%%%%%%%%%%%%%%%%%%%%%%%%%%%%%%%%%%%%
\end{tabular}\\
\vspace{2cm}

\begin{normalsize}October 2026\end{normalsize}

\end{center}

%%%%%%%%%%%%%%%%%%%%%%%%%%%%%%%%%%%%%%%%%%%%%%%%%%%%%%%%%%%%%%%%%%%%%%%%%%%%%%%%%%%%%%%%%%%%%%%%%%
% END OF FILE
%%%%%%%%%%%%%%%%%%%%%%%%%%%%%%%%%%%%%%%%%%%%%%%%%%%%%%%%%%%%%%%%%%%%%%%%%%%%%%%%%%%%%%%%%%%%%%%%%%


\renewcommand{\thepage}{\arabic{page}}
\setcounter{page}{1}

%=============================================================================
\section{Introduction and How to Run}
%=============================================================================
This report describes a small, type-safe arithmetic expression interpreter
in Haskell~\cite{hutton2016programming}, covering Parts A--E of the
assignment. Every module compiles without warnings under GHC~9.6.7 with
\code{-Wall}. Only \code{cabal test} needs more than GHC (it downloads
QuickCheck).
\begin{lstlisting}[style=plain]
git clone https://github.com/Raviyaa27/functional-programming-project
cd "functional-programming-project/Project Files/haskell"
runghc Main.hs     # demo: samples, simplify, batch report
runghc Tests.hs    # 72 expected-vs-actual unit tests
cabal test         # the unit tests and 8 QuickCheck properties
python ../imperative/imperative_eval.py    # Part D (Python 3.10+)
\end{lstlisting}

%=============================================================================
\section{Design and Implementation}\label{sec:design}
%=============================================================================
\subsection{Language Design (Part A)}
\begin{lstlisting}[style=hs, numbers=none, xleftmargin=0pt, framexleftmargin=0pt,
  caption={The types of Part A (abridged from \code{Expr.hs}, lines 47--72).},
  label=lst:types]
data Expr  = Lit Double | Var String
           | Add Expr Expr | Sub Expr Expr | Mul Expr Expr | Div Expr Expr
           | Let String Expr Expr            -- let x = e1 in e2
           | If BExpr Expr Expr              -- if b then e1 else e2
data BExpr = BLit Bool | Less Expr Expr | Equal Expr Expr
           | And BExpr BExpr | Or BExpr BExpr | Not BExpr
type Env   = [(String, Double)]
\end{lstlisting}
\code{Expr} is a sum type with eight alternatives, each carrying exactly the
fields it needs: besides the six required constructors, \code{Let} adds a
local binding and \code{If} a choice. Conditions have their own type,
\code{BExpr}, so numbers and conditions cannot be mixed. In an \code{Env},
\code{lookup} finds the first match, so new bindings shadow old ones.

\subsection{Evaluation (Part B)}
\lstinputlisting[style=hs, firstline=84, lastline=102, firstnumber=84,
  caption={The evaluator (\code{Expr.hs}, lines 84--102).},
  label=lst:eval]{../haskell/Expr.hs}
\code{eval} has one equation per constructor and recurses on the
sub-expressions. Errors are values of type \code{Either String Double}: the
arithmetic cases use applicative style and \code{Div}, \code{Let} and
\code{If} use \code{do}-notation in the \code{Either}
monad~\cite{wadler1995monads}, so the first \code{Left} stops evaluation and
becomes the result. An undefined variable (line~86) and division by zero
(lines 91--96) are reported this way, never as exceptions. The mutually
recursive \code{evalB} (lines 106--116) evaluates conditions with
short-circuiting \code{And}/\code{Or}, and \code{If} evaluates only the chosen
branch.

\subsection{Simplification and Higher-Order Functions (Part C)}
\code{simplify = rewrite . descend} (lines 133--168) simplifies the children
first and then applies the first matching rule at the root: nine identities
($x+0$, $0+x$, $x-0$, $x-x$, $x \cdot 1$, $1 \cdot x$, $x \cdot 0$,
$0 \cdot x$, $x/1$), constant folding (never of $n/0$, which \code{eval} must
report), \code{if} with a constant condition, and removal of an unused
\code{let}; \code{simplifyB} adds double negation and constant comparisons.
Being bottom-up, one pass reaches a normal form, e.g.\
\code{x * 1 + 0}~$\to$~\code{x}.
\lstinputlisting[style=hs, numbers=none, xleftmargin=0pt, framexleftmargin=0pt,
  linerange={204-205,209-209,221-222,224-227},
  caption={The list functions (excerpt of \code{Expr.hs}, lines 204--227).},
  label=lst:hof]{../haskell/Expr.hs}
\code{eval env} is a partially applied (curried) function
\code{Expr -> Either String Double} that \code{map} applies to each expression
(line~205). \code{evalBatch} composes it with \code{rights}, \code{batchReport}
counts successes and failures in one \code{foldr} (6 and 3 on the samples),
and \code{undefinedVars} uses \code{filter} with the section
\code{(\textasciigrave notElem\textasciigrave{} bound)} (line~257) to find unbound variables statically.

%=============================================================================
\section{Sample Evaluations and Testing}\label{sec:samples}
%=============================================================================
Table~\ref{tab:samples} compares hand-computed results for the nine samples
($x=5$, $y=2$, $z=0$) with the output of \code{Main.hs}. \code{Tests.hs} adds
72 unit tests and \code{Props.hs} eight QuickCheck
properties~\cite{claessen2000quickcheck} on 1000 random expressions each; all
pass.

\begin{table}[tb]
\centering
\caption{Sample evaluations: expected vs.\ actual output of \code{Main.hs}.}
\label{tab:samples}
\scriptsize
\renewcommand{\arraystretch}{1.0}
\begin{tabularx}{\textwidth}{@{}c>{\raggedright\arraybackslash\ttfamily}X
    >{\raggedright\arraybackslash\ttfamily}p{3.5cm}
    >{\raggedright\arraybackslash\ttfamily}p{3.5cm}c@{}}
\toprule
\textbf{\#} & \textnormal{\textbf{Expression}} & \textnormal{\textbf{Expected}}
  & \textnormal{\textbf{Actual}} & \textbf{OK} \\
\midrule
1 & (x + 3) * y & 16 & 16 & \checkmark \\
2 & x / y & 2.5 & 2.5 & \checkmark \\
3 & let a = x * 2 in a + 1 & 11 & 11 & \checkmark \\
4 & let x = 1 in (let x = 10 in x) + x & 11 & 11 & \checkmark \\
5 & if x < 10 \&\& not (y == 0) then x / y else 0 & 2.5 & 2.5 & \checkmark \\
6 & if z == 0 then 0 else x / z & 0 & 0 & \checkmark \\
7 & x / (y - 2) & error: division by zero: x / (y - 2)
  & error: division by zero: x / (y - 2) & \checkmark \\
8 & x + w & error: undefined variable: w & error: undefined variable: w
  & \checkmark \\
9 & let a = 1 / z in a + 1 & error: division by zero: 1 / z
  & error: division by zero: 1 / z & \checkmark \\
\bottomrule
\end{tabularx}
\end{table}

%=============================================================================
\section{Comparative Analysis (Part D)}\label{sec:compare}
%=============================================================================
Listing~\ref{lst:python} is the imperative version: a \code{match} on string
tags, a mutable \code{result}, a shared dictionary and exceptions.

\begin{table}[tb]
\centering
\caption{The Haskell and Python versions compared.}
\label{tab:compare}
\footnotesize
\renewcommand{\arraystretch}{1.0}
\begin{tabularx}{\textwidth}{@{}>{\raggedright\arraybackslash}p{2.5cm}
    >{\raggedright\arraybackslash}X>{\raggedright\arraybackslash}X@{}}
\toprule
 & \textbf{Haskell (\code{Expr.hs})} & \textbf{Python (\code{imperative\_eval.py})} \\
\midrule
Error representation & \code{Left msg}, a value; the type shows that failure
  is possible & \code{EvalError} exception; not visible in the signature \\
Error propagation & \code{<*>} and \code{do} stop at the first \code{Left};
  callers must unwrap the result & Stack unwinding to whichever \code{try}
  encloses the call \\
Mutable state & None: \code{let} builds a new list \mbox{\code{(x, v) : env}} &
  \code{result} reassigned; \code{env} overwritten, restored in
  \code{finally} \\
Misspelt or missing case & Compile error (\code{Mull} not in scope);
  \code{-Wincomplete-patterns} warning & Found only if that
  branch runs (\code{case \_}) \\
Mixing numbers and conditions & Type error: \code{If} needs a \code{BExpr},
  \code{Add} needs \code{Expr}s & Accepted: any value is ``truthy'', and
  \code{True + 3} gives \code{4} \\
\bottomrule
\end{tabularx}
\end{table}

\textbf{Errors.} In Haskell, failure is part of the type of \code{eval}, so
a caller must handle the \code{Left} before using the \code{Double}, and
\code{<*>} and \code{do} propagate it at no cost. In Python, failure is
invisible in the signature and a forgotten \code{except} crashes the
program.
\textbf{State.} Python reassigns \code{result} in every case and overwrites the
shared \code{env} in \code{let}; without the \code{finally}, an error in a
\code{let} body would leave the binding behind. Haskell has no mutable state.
\textbf{Bugs ruled out.} \code{imperative\_eval.py} runs five programs that
GHC rejects before they run (the last rows of Table~\ref{tab:compare}, plus a
string literal that crashes with \code{TypeError}). Division by zero and
unbound variables depend on run-time values, so neither version excludes them
by types; Haskell returns them as values instead of throwing them.

\lstinputlisting[style=py, linerange={16-47,61-63},
  caption={Imperative evaluator (\code{imperative\_eval.py}, lines 16--63;
  the condition cases are omitted).},
  label=lst:python]{../imperative/imperative_eval.py}

%=============================================================================
\section{Reflection (Part E)}\label{sec:reflect}
%=============================================================================
\textbf{Pure functions.} A pure function's result depends only on its
arguments, and it has no side effects. The type of \code{eval} (line~84) has
no \code{IO}, so it cannot print, read input or change any variable; the same
environment and expression always give the same result, and likewise for
\code{simplify}. This made testing simple: a unit test just compares values,
and QuickCheck can compare \code{eval env e} with
\code{eval env (simplify e)} thousands of times. Functions are also values:
\code{eval env} is passed to \code{map} (line~205)~\cite{hughes1989why}.

\textbf{Immutability.} Nothing is updated in place. Line~99,
\code{eval ((x, v) : env) body}, does not modify \code{env}: it builds a new
list whose tail \emph{is} the old one (shared, not copied). When the
\code{let} ends, the caller still holds its original \code{env}, so the
binding disappears with no clean-up (sample~4) and Python's
\code{try}/\code{finally} has no counterpart. \code{simplify} likewise builds a
new tree (lines 140--145) and leaves its input intact, so \code{Main.hs}
prints each expression before and after simplification from the same value.

\textbf{Why invalid expressions cannot be constructed.} Each constructor fixes
the number and types of its fields: \code{Add} takes exactly two \code{Expr}s,
and \code{If} takes a \code{BExpr} and two \code{Expr}s. An operator with a
missing operand, or a misspelt one such as \code{Mull}, is therefore not an
\code{Expr} at all, and keeping conditions in a separate type turns Python's
silent bugs into compile errors (\code{compile\_errors.txt}):
\begin{lstlisting}[style=plain]
bad1 = If (Lit 1) (Lit 2) (Lit 3)
  * Couldn't match expected type `BExpr' with actual type `Expr'
\end{lstlisting}
\code{Add (Less (Lit 1) (Lit 2)) (Lit 3)} fails likewise with the types
swapped, and a function that forgets a constructor gets a warning naming it
(\code{Patterns of type \textasciigrave Expr\textquotesingle{} not matched: If \_ \_ \_}). Errors that depend
on values (an unbound name, a zero divisor) are beyond the types; \code{eval}
returns them as \code{Left}, and \code{undefinedVars} finds unbound names
before evaluation (QuickCheck confirms it is sound and complete).
\code{TypedExpr.hs} goes further with a GADT indexed by the result type, so a
single expression type still rejects ill-typed terms at compile time.

\textbf{What testing revealed.} QuickCheck refuted the claim that
\code{simplify} never changes whether evaluation succeeds by finding
\code{w - w}: it fails with \code{undefined variable: w} but simplifies to
\code{0}, because rules that discard a sub-expression ($x \cdot 0$, $x - x$,
an unused \code{let}) discard its error too. What holds, on 1000 random
expressions, is that every \emph{successful} result is preserved; because
failure is a value, this was visible and testable.

% References
\bibliographystyle{ieeetr}
\bibliography{bibliography}

\end{document}
